\chapter{C$^*$-algebras, states, and representations}\label{ch:cstar_states}

The algebraic formulation of quantum theory: observables as a C$^*$-algebra, states as positive functionals, the GNS construction that turns a state into a Hilbert-space representation, the inequivalent representations that appear in infinite systems, and the thermodynamic limit.

\section{\cstar-algebras}\label{ch:cstar}

\begin{physmotiv}
What do we know about the observables of a system before we choose a state or a Hilbert space? That they can be added, multiplied by complex numbers, multiplied together (non-commutatively), conjugated ($A\mapsto A\ad$), and that they have a norm (a measure of the largest possible measured value). A \cstar-algebra is exactly this: nothing more, nothing less. For a system with finitely many degrees of freedom the algebra is $\BH$; for a spin chain or a field one has to construct it directly from the local commutation relations, since no preferred $\HH$ exists (\cref{ch:why_algebras}).
\end{physmotiv}

\subsection{Definition and the basic theorem}

\begin{definition}
A \emph{\cstar-algebra} is a complex algebra $\A$ with an involution $a\mapsto a\ad$ and a norm $\norm\cdot$ such that $\A$ is complete, $\norm{ab}\le\norm a\norm b$, and
\begin{equation}
\norm{a\ad a}=\norm a^2\qquad(\text{the \cstar-identity}).
\end{equation}
We take $\A$ unital (with unit $\id$) unless stated otherwise.
\end{definition}

The \cstar-identity ties the norm to the algebra. It implies that the norm is determined by the algebraic structure: $\norm a^2$ is the spectral radius of $a\ad a$. Hence an algebra has \emph{at most one} \cstar-norm, and every $^*$-homomorphism between \cstar-algebras is automatically norm-decreasing.

\begin{theorem}[Gelfand--Naimark]
Every \cstar-algebra is isometrically $^*$-isomorphic to a norm-closed $^*$-subalgebra of $\Bnd{\HH}$ for some Hilbert space $\HH$.
\end{theorem}
Thus abstract \cstar-algebras and norm-closed algebras of operators are the same thing, but the Hilbert space is not unique. The construction of such a $\HH$ from a state is the GNS construction (\cref{ch:gns}).

\subsubsection*{Spectrum, positivity, functional calculus}
An element $a$ is self-adjoint if $a=a\ad$. The \emph{spectrum} $\spec(a)=\{\lambda\in\C:a-\lambda\id\text{ not invertible}\}$ is a non-empty compact set, real for self-adjoint $a$, and \emph{independent of which algebra} (or representation) it is computed in: this is crucial since it means $\spec(a)$ is intrinsic to $\A$. For self-adjoint $a$ one has the \emph{continuous functional calculus}: $f(a)\in\A$ for every continuous function $f$ on $\spec(a)$. An element is \emph{positive} ($a\ge0$) if it is self-adjoint with $\spec(a)\subset[0,\infty)$; equivalently $a=b\ad b$ for some $b$. Positivity gives the order structure on observables used to define states (\cref{ch:states}).

\begin{whatwrong}
Only \emph{continuous} functions of $a$ lie in the \cstar-algebra. Spectral projections $\chi_\Omega(a)$ (functions with jumps), in general do \emph{not}: take $\A=C[0,1]$ and $a(x)=x$; then $\chi_{[0,1/2]}(a)$ is discontinuous and not in $C[0,1]$. They do lie in the von Neumann algebra obtained by weak closure (\cref{ch:doublecomm}). This is a reason one passes from \cstar\ to vN algebras: yes/no questions need to be elements.
\end{whatwrong}

\subsection{Examples}

\begin{enumerate}
\item \textbf{Matrices} $M_n(\C)=\Bnd{\C^n}$. The universal finite-dimensional example; every finite-dimensional \cstar-algebra is a direct sum $\bigoplus_kM_{n_k}(\C)$.
\item \textbf{Compact operators} $\mathcal K(\HH)$ (non-unital) and $\BH$.
\item \textbf{Commutative algebras} $C(X)$, continuous complex functions on a compact Hausdorff space $X$ with the sup-norm. By the Gelfand--Naimark \emph{duality} (\cref{ch:gelfand}) every commutative unital \cstar-algebra has this form with $X$ the space of characters. Commutative \cstar-algebras are classical phase spaces.
\item \textbf{Weyl (CCR) algebra} $\mathcal W(\mathcal S,\sigma)$ of a symplectic space: the \cstar-algebra generated by unitaries $W(f)$ with $W(f)W(g)=\ee^{\frac\ii2\sigma(f,g)}W(f+g)$ (\cref{ch:weyl}). It has a unique \cstar-norm (Slawny), so it is a well-defined object independent of representation. For a free scalar field of mass $m$ in Minkowski space, $\sigma$ is the Pauli--Jordan commutator and the algebra of a region $R$ is the subalgebra $\mathcal W(\mathcal S_R)$ generated by $W(f)$ with $f$ supported in $R$.
\item \textbf{CAR algebra} of fermions: generated by $a(f)$, $f\in\mathcal K$ a Hilbert space, with $\{a(f),a(g)\ad\}=\inner fg\id$, $\{a(f),a(g)\}=0$. Since $a(f)$ is automatically bounded ($\norm{a(f)}=\norm f$), no Weyl trick is needed.
\item \textbf{The infinite spin chain (UHF algebra).} Let $\A_N=\bigotimes_{j=1}^NM_2(\C)=M_{2^N}(\C)$, embedded as $\A_N\hookrightarrow\A_{N+1}$ via $a\mapsto a\otimes\id$. The norm-closure of $\bigcup_N\A_N$ is the \cstar-algebra $\A=M_{2^\infty}$ of the infinite chain. It is \emph{simple} (no non-trivial closed two-sided ideals) and unique: one algebra, an infinite family of different representations (\cref{ch:inequiv}), as in \cref{ch:why_algebras}. It carries a unique \emph{tracial} state $\tau(ab)=\tau(ba)$, the infinite-temperature state, which we will meet again as the origin of the hyperfinite type II$_1$ factor (\cref{ch:classification}).
\end{enumerate}

\begin{keyidea}
A \cstar-algebra is the representation-independent algebra of observables. Matrices, classical phase spaces (commutative case), CCR/CAR algebras, and spin-chain algebras are all \cstar-algebras. Continuity of the norm means only \emph{continuous} functions of observables are in it.
\end{keyidea}

\subsection{Locality and the quasi-local picture}

Algebras of \emph{local observables}, those supported in a finite region $\Lambda$, are $\A(\Lambda)$ for each finite $\Lambda$ and one has isotony $\Lambda_1\subset\Lambda_2\Rightarrow\A(\Lambda_1)\subset\A(\Lambda_2)$, and commutation $[\A(\Lambda_1),\A(\Lambda_2)]=0$ if $\Lambda_1\cap\Lambda_2=\emptyset$ (for bosons/spins; graded for fermions). The \emph{quasi-local algebra} is the norm closure $\A=\overline{\bigcup_\Lambda\A(\Lambda)}$. This is the infinite spin chain example and the framework of \cref{ch:quasilocal}. In relativistic QFT we replace $\Lambda$ by spacetime regions (\cref{ch:haagkastler}).

\subsection*{Exercises}
\begin{exercise}\label{ex:cstar:1}
Show from the \cstar-identity that $\norm{a\ad}=\norm a$ and that for self-adjoint $a$, $\norm{a^{2^k}}=\norm a^{2^k}$, hence $\norm a=\lim\norm{a^n}^{1/n}$ (spectral radius). Conclude the norm is determined by the algebra.
\end{exercise}
\begin{exercise}\label{ex:cstar:2}
Show that $\A_N=M_{2^N}$ and the map $a\mapsto a\otimes\id$ is an isometry, so that the inductive limit norm is well defined. Show that $\sigma^z_j$ and $\sigma^x_k$ commute for $j\ne k$. Is $M_N=\frac1N\sum_{j\le N}\sigma_j^z$ in $\A$? Does the sequence converge in norm? (Compare \cref{ch:topologies}.)
\end{exercise}
\begin{exercise}\label{ex:cstar:3}
For $\A=C[0,1]$, find the spectrum of $a(x)=x$, compute $\norm a$, and show $\chi_{[0,1/2]}(a)\notin\A$. Show that every closed two-sided ideal of $C(X)$ is of the form $\{f:f|_Y=0\}$ for a closed $Y\subset X$.
\end{exercise}
\begin{exercise}\label{ex:cstar:4}
Show that for fermionic modes $a(f)$ one has $\norm{a(f)}=\norm f$ using $a(f)a(f)\ad+a(f)\ad a(f)=\norm f^2\id$. Explain why one can treat fermions without exponentiating.
\end{exercise}

\section{States}\label{ch:states}

\begin{physmotiv}
An experiment produces a list of expectation values: for each observable $a$ a number $\omega(a)$. The map $a\mapsto\omega(a)$ is linear, real on self-adjoint elements, non-negative on positive ones (probabilities are non-negative), and sends the identity to $1$. These are the \emph{only} properties used. A state is such a map: no vector, no density matrix, no Hilbert space. In finite dimensions the map is $\omega(a)=\Tr\rho a$; in general it need not be.
\end{physmotiv}

\subsection{Definition}

\begin{definition}
A \emph{state} on a unital \cstar-algebra $\A$ is a linear functional $\omega:\A\to\C$ which is \emph{positive}, $\omega(a\ad a)\ge0$ for all $a$, and \emph{normalized}, $\omega(\id)=1$.
\end{definition}

Basic consequences (all easy):
\begin{itemize}
\item $\omega(a\ad)=\overline{\omega(a)}$, and positivity implies automatic norm-continuity: $|\omega(a)|\le\norm a$.
\item \textbf{Cauchy--Schwarz:} $|\omega(b\ad a)|^2\le\omega(b\ad b)\,\omega(a\ad a)$. This is the algebraic origin of the inner product in the GNS construction.
\item The set $\mathcal S(\A)$ of states is convex and compact in the weak-$^*$ topology (convergence of $\omega_\alpha(a)\to\omega(a)$ for each $a$). The latter follows from Banach--Alaoglu.
\end{itemize}

\subsection{Pure vs.\ mixed}

A state is \emph{mixed} if it is a non-trivial convex combination, $\omega=\lambda\omega_1+(1-\lambda)\omega_2$ with $0<\lambda<1$ and $\omega_1\ne\omega_2$, and \emph{pure} otherwise (an extreme point of $\mathcal S(\A)$).

\begin{center}
\renewcommand{\arraystretch}{1.3}
\resizebox{\ifdim\width>\linewidth\linewidth\else\width\fi}{!}{%
\begin{tabular}{@{}lll@{}}
\toprule
Algebra & States & Pure states\\
\midrule
$M_n(\C)$ & density matrices $\rho$, $\omega(a)=\Tr\rho a$ & $\rho=\ketbra\psi\psi$\\
$C(X)$ (classical) & probability measures on $X$ (Riesz) & Dirac measures $\delta_x$ (points of $X$)\\
$\BH$, $\dim\HH=\infty$ & normal (density matrix) \emph{and} singular states & vector states; also singular ones\\
\bottomrule
\end{tabular}}
\end{center}

The commutative case shows how classical probability fits: a classical state is a probability measure, and its pure states are deterministic points. In the non-commutative case pure states are far more constrained: for $M_2(\C)$ the state space is the Bloch ball, pure states are the sphere, and a mixed state has \emph{many} different decompositions into pure states, whereas in $C(X)$ each state has an essentially unique decomposition into points.

\begin{whatwrong}
For $\BH$ with $\dim\HH=\infty$ a state is \emph{not} always a density matrix. By the Hahn--Banach theorem there are states that vanish on all compact operators: they assign probability $0$ to every finite-rank projection, so they look like particles ``at infinity''. They are \emph{singular} (not normal). Their GNS representations are inequivalent to the identity representation. In the von Neumann algebra framework (\cref{ch:traces}) we will restrict attention to \emph{normal} states, which are the continuous ones, i.e.\ those given by density matrices if the algebra is $\BH$.
\end{whatwrong}

\subsection{Examples}

\begin{enumerate}
\item \textbf{Product states on the spin chain.} For $\A=M_{2^\infty}$ and $\theta\in[0,\pi/2]$, $\omega_\theta=\bigotimes_j\omega^{(j)}_\theta$ with $\omega^{(j)}_\theta(a)=\bra{\psi_\theta}a\ket{\psi_\theta}$. Then $\omega_\theta(\sigma^z_j)=\cos2\theta$. Pure states, all inequivalent for different $\theta$ (\cref{ch:inequiv}).
\item \textbf{Infinite-temperature state.} $\tau=\bigotimes_j\tfrac12\Tr$, the unique tracial state: $\tau(ab)=\tau(ba)$.
\item \textbf{Gibbs states and thermodynamic limits.} For a finite lattice $\Lambda$, $\omega_\Lambda(a)=\Tr(\ee^{-\beta H_\Lambda}a)/Z_\Lambda$. The weak-$^*$ limits of $\omega_\Lambda$ as $\Lambda\uparrow\Z^d$, if they exist, define a state on the quasi-local algebra even though $\ee^{-\beta H}$ makes no sense in infinite volume (\cref{ch:kms,ch:quasilocal}). Compactness of $\mathcal S(\A)$ ensures that limit points exist; non-uniqueness of the limit signals a phase transition.
\item \textbf{Vacuum of a free field.} On the Weyl algebra the Fock state is $\omega(W(f))=\ee^{-\frac12\norm{Kf}^2}$ where $K$ maps the test data to the one-particle Hilbert space; it is a \emph{quasi-free} state, completely determined by its two-point function.
\end{enumerate}

\subsection{Vector states, normal states, and what a state ``knows''}

Given any representation $\pi:\A\to\BH$ and unit vector $\xi\in\HH$, $\omega_\xi(a)=\inner\xi{\pi(a)\xi}$ is a state (a \emph{vector state}). The word ``pure'' refers to the algebra on which the state is considered, not to the vector: in the algebra of the right Rindler wedge the vacuum state is a \emph{mixed} (indeed thermal) state although $\Omega$ is a unit vector in $\HH$. A pure state restricted to a subalgebra is typically mixed, which is the algebraic form of entanglement: in $M_2\otimes M_2$ a pure entangled state restricts to a mixed state on $M_2\otimes\id$.

\begin{keyidea}
A state is a positive normalized linear functional. It is the only operational information about a quantum system. The notion of a vector in a Hilbert space arises from it by GNS (next section).
\end{keyidea}

\subsection*{Exercises}
\begin{exercise}\label{ex:states:1}
Derive the Cauchy--Schwarz inequality for a state from positivity of $\omega\big((a+\lambda b)\ad(a+\lambda b)\big)$ for all $\lambda\in\C$.
\end{exercise}
\begin{exercise}\label{ex:states:2}
Characterize all states on $M_2(\C)$ as a Bloch ball. Show that a mixed state has infinitely many decompositions into pure states. Contrast with the commutative algebra $\C^2=C(\{0,1\})$ (a classical bit), where every state has a unique decomposition into pure states.
\end{exercise}
\begin{exercise}\label{ex:states:3}
Let $\ket\psi=\cos\alpha\ket{\uparrow\uparrow}+\sin\alpha\ket{\downarrow\downarrow}$ and $\A=M_2\otimes M_2$. Show that the state $\omega_\psi$ is pure on $\A$ but its restriction to $M_2\otimes\id$ is mixed unless $\alpha\in\{0,\pi/2\}$, and compute the entropy.
\end{exercise}
\begin{exercise}\label{ex:states:4}
Show that the infinite-temperature state $\tau$ on the spin chain is tracial and that it is the unique tracial state. (Hint: any tracial state restricted to $M_{2^N}$ must be $\Tr/2^N$.)
\end{exercise}
\begin{exercise}\label{ex:states:5}
For $\A=\BH$ with $\HH$ infinite dimensional, show the set of normal states is a convex proper subset of $\mathcal S(\A)$. Prove existence of a singular state by considering the quotient $\BH/\mathcal K(\HH)$ and a Hahn--Banach extension.
\end{exercise}

\section{The GNS construction}\label{ch:gns}

\begin{physmotiv}
If observables and states are the primary objects, where do Hilbert spaces and state vectors come from? The answer is a beautiful and completely explicit construction, due to Gelfand, Naimark and Segal. Given an algebra $\A$ and a state $\omega$, one builds a Hilbert space $\HH_\omega$, a representation $\pi_\omega$ of $\A$ on it, and a vector $\Omega_\omega$ such that $\omega(a)=\inner{\Omega_\omega}{\pi_\omega(a)\Omega_\omega}$. It is the rigorous version of ``act on the vacuum with operators to produce all states'', and it shows \emph{how} different states give different Hilbert spaces. The cyclic vector $\Omega$ is exactly the object that will give us the modular operator (\cref{ch:tomita}).
\end{physmotiv}

\subsection{The finite-dimensional example first}

Let $\A=M_n(\C)$ and $\omega(a)=\Tr(\rho a)$ with $\rho>0$ (faithful). Try to use $\A$ itself as the Hilbert space, with the inner product
\begin{equation}
\inner ab_\omega:=\omega(a\ad b)=\Tr(\rho\,a\ad b).
\end{equation}
This is positive definite since $\omega(a\ad a)=\Tr(\rho^{1/2}a\ad a\rho^{1/2})\ge0$, with equality iff $a\rho^{1/2}=0$, i.e.\ $a=0$ since $\rho$ is invertible. Define $\pi(a)b=ab$ (left multiplication) and $\Omega=\id\in\A$. Then $\inner\Omega{\pi(a)\Omega}=\omega(a)$. Now use the map $b\mapsto b\rho^{1/2}$ to identify $\A$ with the Hilbert--Schmidt operators $\mathrm{HS}\cong\C^n\otimes\overline{\C^n}$: $\inner ab_\omega=\Tr\big((a\rho^{1/2})\ad(b\rho^{1/2})\big)$. In this picture
\begin{equation}
\Omega=\rho^{1/2}=\sum_i\sqrt{p_i}\,\ket i\otimes\ket{\bar i},
\end{equation}
which is the \emph{purification} of $\rho$ (Schmidt decomposition with weights $p_i$), $\pi(a)=a\otimes\id$, and the commutant is $\pi(\A)'=\id\otimes\overline{M_n}$ (right multiplication). Therefore:

\begin{keyidea}
The GNS representation of a mixed state of $M_n$ is the purification: $\HH_\omega=\HH\otimes\overline\HH$, $\Omega=\sum\sqrt{p_i}\ket{i}\ket{\bar i}$. The ``doubled'' system is not an artifact: it is the commutant $\pi(\A)'$, and it will become the algebra of the complementary region in QFT (Rindler left wedge; thermofield double).
\end{keyidea}

\subsection{The general construction}

Let $\A$ be a unital \cstar-algebra and $\omega$ a state.
\begin{enumerate}
\item By Cauchy--Schwarz, $N_\omega=\{a\in\A:\omega(a\ad a)=0\}$ is a \emph{left ideal}; it is the set of ``unobservably small'' vectors. (Indeed $\omega(b\ad a)=0$ for all $b$ and $a\in N_\omega$.)
\item On the quotient $\A/N_\omega$ define $\inner{[a]}{[b]}=\omega(a\ad b)$. This is a positive definite inner product. Let $\HH_\omega$ be the completion.
\item Define $\pi_\omega(a)[b]=[ab]$. This is well defined because $N_\omega$ is a left ideal, and extends to bounded operators with $\norm{\pi_\omega(a)}\le\norm a$.
\item Set $\Omega_\omega=[\id]$.
\end{enumerate}

\begin{theorem}[GNS]\label{thm:GNS}
$(\HH_\omega,\pi_\omega,\Omega_\omega)$ is a representation of $\A$ with a \emph{cyclic} unit vector $\Omega_\omega$ (i.e.\ $\pi_\omega(\A)\Omega_\omega$ is dense in $\HH_\omega$) satisfying $\omega(a)=\inner{\Omega_\omega}{\pi_\omega(a)\Omega_\omega}$. It is unique up to unitary equivalence: any cyclic representation $(\HH,\pi,\Omega)$ with $\omega=\inner\Omega{\pi(\cdot)\Omega}$ is unitarily equivalent to it by a unitary mapping $\Omega\mapsto\Omega_\omega$.
\end{theorem}

\begin{proofsketch}
Positivity of $\omega$ gives positivity of the inner product, and the \cstar-algebra properties (and boundedness of left multiplication, which follows from $a\ad b\ad ba\le\norm b^2a\ad a$ in the operator order) give boundedness of $\pi_\omega(b)$. Uniqueness: define $U\pi_\omega(a)\Omega_\omega=\pi(a)\Omega$. It preserves inner products since both equal $\omega(a\ad b)$; so it extends to a unitary on the closures, which are everything by cyclicity.
\end{proofsketch}

\begin{whatwrong}
$\Omega_\omega$ is \emph{cyclic} by construction but need not be \emph{separating} for $\pi_\omega(\A)''$ (i.e.\ it may be annihilated by non-zero operators). For a pure state on $M_n$, $n\ge2$, $\Omega=\psi$, $\pi(\A)=M_n$ and any $a$ with $a\psi=0$ is nonzero: $\psi$ is cyclic but not separating. The state is not faithful. For the Rindler wedge and the vacuum $\Omega$, both properties hold (\cref{ch:rs_bw}); this is what makes modular theory possible. Cyclic and separating vectors are therefore precisely those of ``full-rank'' states.
\end{whatwrong}

\subsection{Pure states and irreducibility}

\begin{theorem}
$\omega$ is pure if and only if $\pi_\omega$ is irreducible (no non-trivial closed invariant subspace; equivalently $\pi_\omega(\A)'=\C\id$).
\end{theorem}
\emph{Idea.} Any positive $T\in\pi_\omega(\A)'$ with $0\le T\le\id$ defines a positive functional $\omega_T(a)=\inner{\Omega}{T\pi_\omega(a)\Omega}$ dominated by $\omega$; conversely every functional $\le\omega$ arises this way. If $\omega=\lambda\omega_1+(1-\lambda)\omega_2$ then $\lambda\omega_1\le\omega$ gives a $T$ not proportional to $\id$, and conversely. Thus mixed states have non-trivial commutant (they carry ``classical'' information, the $T$), and pure states do not.

Classic consequence (Schur-type): if $\pi_\omega$ is irreducible then the von Neumann algebra $\pi_\omega(\A)''=\BH_\omega$, all bounded operators. Every pure state on a \cstar-algebra is a vector state of a full $\Bnd{\HH_\omega}$, matching textbook QM, but \emph{different pure states give different, perhaps inequivalent, $\HH_\omega$}.

\subsection{Worked examples}

\subsubsection{Fock representation from the Weyl algebra}
On $\A=\mathcal W(\mathcal S,\sigma)$ take a one-particle structure, i.e.\ a real-linear map $K:\mathcal S\to\mathfrak h$ into a complex Hilbert space with $\sigma(f,g)=2\,\mathrm{Im}\inner{Kf}{Kg}$ (up to a sign convention) and $K\mathcal S+\ii K\mathcal S$ dense. The quasi-free state $\omega(W(f))=\ee^{-\frac12\norm{Kf}^2}$ is positive, and its GNS representation is the \emph{Fock representation} on the symmetric Fock space $\mathcal F(\mathfrak h)$, with $\Omega$ the Fock vacuum and $W(f)=\exp\ii\big(a(Kf)+a(Kf)\ad\big)$. The vector $\Omega$ is cyclic because the $W(f)\Omega$ are coherent states, which span a dense set. The choice of $K$ is a choice of vacuum (positive-frequency structure). Different choices give different states; if they are not related by a unitarily implementable Bogoliubov transformation, they give inequivalent representations (\cref{ch:weyl,ch:inequiv}).

\subsubsection{The vacuum of a free field}
For the free scalar field in Minkowski space, taking $K$ to be the positive-frequency part of the solution restricts the Fock construction to the Poincar\'e-invariant vacuum state $\omega_0$ of the quasi-local algebra, and the GNS representation realizes the Wightman fields as operators on $\HH_0=\mathcal F(L^2(\R^d,\dd^dk/2\omega_k))$: the vacuum representation used in textbook QFT. All the standard QFT objects, the two-point function $\inner\Omega{\phi(x)\phi(y)\Omega}$, one-particle states, and so on, can be reconstructed from the state alone.

\subsubsection{Thermal states}
The state $\omega_\beta$ on the Weyl algebra with $\omega_\beta(W(f))=\exp\big[-\tfrac12\inner{Kf}{\coth(\beta\omega/2)Kf}\big]$ (the Gibbs state of the free field in the thermodynamic limit, with $\omega$ now denoting the one-particle energy operator) has GNS space $\mathcal F(\mathfrak h)\otimes\mathcal F(\overline{\mathfrak h})$, the \emph{thermofield double}\index{thermofield double} Hilbert space, with $\Omega_\beta$ the thermofield-double state. We will see this again in \cref{ch:kms,ch:tomita}; compare with $\HH\otimes\overline\HH$ in the finite-dimensional example.

\subsection*{Exercises}
\begin{exercise}\label{ex:gns:1}
Carry out the GNS construction for $\A=M_2(\C)$ and (a) a pure state $\omega=\bra\psi\cdot\ket\psi$, (b) $\rho=\mathrm{diag}(p,1-p)$, $0<p<1$. Find $N_\omega$, $\HH_\omega$ (dimension), $\pi_\omega(\A)'$, and decide whether $\Omega$ is separating.
\end{exercise}
\begin{exercise}\label{ex:gns:2}
For $\A=C(X)$ and the state $\omega(f)=\int f\,\dd\mu$, show that $\HH_\omega=L^2(X,\mu)$, $\pi_\omega$ is multiplication, and $\Omega=1$. Show $\pi_\omega(\A)'=\Linf(X,\mu)$ (maximal abelian), and that $\pi_\omega$ is irreducible iff $\mu$ is a point mass.
\end{exercise}
\begin{exercise}\label{ex:gns:3}
For the spin chain $M_{2^\infty}$ and the tracial state $\tau$, show that $\HH_\tau$ is the closure of $\A$ in the norm $\tau(a\ad a)^{1/2}$. Show that right multiplication gives $\pi_\tau(\A)'$, which is isomorphic to $\pi_\tau(\A)''$ (to be explained in \cref{ch:classification}); the weak closure $\pi_\tau(\A)''$ is the hyperfinite II$_1$ factor.
\end{exercise}
\begin{exercise}\label{ex:gns:4}
In the finite-dimensional example, show that $\Omega=\rho^{1/2}$ is cyclic and, for $\rho$ invertible, separating for $\pi(\A)=M_n\otimes\id$. Define the antilinear map $S\,\pi(a)\Omega=\pi(a\ad)\Omega$; in the picture $\HH\otimes\overline\HH$ show that $S=J\,(\rho^{1/2}\otimes\rho^{-1/2})$ where $J$ is the antiunitary swapping the two factors, and compute $S\ad S=\rho\otimes\rho^{-1}$. (This is the modular operator of \cref{ch:tomita} in the simplest case.)
\end{exercise}

\section{Inequivalent representations}\label{ch:inequiv}

\begin{physmotiv}
Two ferromagnets, one magnetized up and one down, are described by the \emph{same} algebra of local observables but not by the same Hilbert space: no finite number of local operations turns one into the other. Phases of matter, superselection sectors, spontaneous symmetry breaking and thermal states in infinite volume are all instances of one mathematical fact: a \cstar-algebra has many inequivalent representations, and GNS (\cref{ch:gns}) shows that they correspond to different states. This section makes that precise and gives practical criteria.
\end{physmotiv}

\subsection{Equivalence, disjointness, quasi-equivalence}

Two representations $\pi_1,\pi_2$ of $\A$ on $\HH_1,\HH_2$ are \emph{unitarily equivalent} if $\pi_2(a)=U\pi_1(a)U\ad$ for a unitary $U:\HH_1\to\HH_2$. They are \emph{disjoint} if no non-zero subrepresentation of one is equivalent to a subrepresentation of the other, and \emph{quasi-equivalent}\index{quasi-equivalence} if each is equivalent to a multiple (direct sum of copies) of a subrepresentation of the other; this holds iff the map $\pi_1(a)\mapsto\pi_2(a)$ extends to an isomorphism of the von Neumann algebras $\pi_1(\A)''\cong\pi_2(\A)''$.

For a representation $\pi$ the \emph{folium}\index{folium} is the set of states given by density matrices in $\pi$ (normal states of $\pi(\A)''$). Two states $\omega_1,\omega_2$ have quasi-equivalent GNS representations iff they lie in the same folium. Physically:

\begin{center}
\emph{States in the same folium differ by ``finitely many'' local excitations: they are experimentally indistinguishable by local operations on a finite region, up to an arbitrarily small error, and they describe the same phase.}
\end{center}
States in different (disjoint) folia describe different phases or sectors: they cannot be obtained from each other by any local operations (or normal perturbations). For \emph{pure} states, quasi-equivalence of GNS representations is the same as unitary equivalence.

\subsection{A decisive tool: central macroscopic observables}

Recall from \cref{ch:topologies} that in the product representation $\pi_\theta$ the average magnetization $M_N=\frac1N\sum_{j\le N}\sigma^z_j$ converges strongly to the scalar $m(\theta)\id=\cos2\theta\,\id$. If $\pi_\theta$ and $\pi_{\theta'}$ were unitarily equivalent via $U$, then $U\big(m(\theta)\id\big)U\ad=\lim UM_NU\ad=\lim M_N=m(\theta')\id$, forcing $m(\theta)=m(\theta')$. Hence
\begin{equation}
\theta\ne\theta'\Longrightarrow\pi_\theta\not\simeq\pi_{\theta'},\quad\text{in fact they are disjoint.}
\end{equation}
\emph{General principle:} a strong limit of local observables which is a scalar in one representation and a different scalar in another proves inequivalence. In the weak closure $\pi(\A)''$, such limits are \emph{central} elements; they are the ``order parameters'' that label phases.

\subsection{A computable criterion: product states (Kakutani)}

Let $\omega=\bigotimes_j\omega_j$ and $\omega'=\bigotimes_j\omega'_j$ be infinite product \emph{vector} states of the spin chain, with local vectors $\psi_j,\psi_j'\in\C^2$.

\begin{theorem}[Kakutani]
$\pi_\omega$ and $\pi_{\omega'}$ are unitarily equivalent if $\sum_j\big(1-|\braket{\psi_j}{\psi_j'}|\big)<\infty$ (equivalently $\prod_j|\braket{\psi_j}{\psi'_j}|>0$), and disjoint otherwise.
\end{theorem}
For the states $\psi_\theta$ at every site, $|\braket{\psi_\theta}{\psi_{\theta'}}|=\cos(\theta-\theta')<1$ is constant, so the product $\to0$ and they are disjoint, consistent with the above. The criterion says: two states are equivalent iff they differ \emph{significantly} on only finitely many sites. This is the lattice version of ``a finite number of excitations'' versus ``a finite density of excitations''.

\begin{whatwrong}
A state which differs from another by a small amount at every site, however small, is in a different sector (for fixed non-summable perturbation). This is why thermodynamic limits cannot be treated as small perturbations of ground states, and why exponentiating a perturbation $\ee^{-\beta H_{\rm int}}$ in infinite volume may leave the Hilbert space.
\end{whatwrong}

\subsection{Spontaneous symmetry breaking}

Let $G$ act on $\A$ by automorphisms $\alpha_g$ (for example the spin flip $\sigma^z_j\mapsto-\sigma^z_j$, $\sigma^{x,y}_j\mapsto\sigma^{x,y}_j$ combined with a rotation, which maps $\psi_\theta\mapsto\psi_{\pi/2-\theta}$ and $m\to-m$). A state $\omega$ is \emph{invariant} if $\omega\circ\alpha_g=\omega$. If $\omega$ is invariant then $U(g)\pi_\omega(a)\Omega=\pi_\omega(\alpha_g(a))\Omega$ defines a unitary representation implementing $\alpha_g$ on $\HH_\omega$, leaving $\Omega$ fixed.

\begin{definition}
The symmetry $\alpha_g$ is \emph{spontaneously broken} in the representation $\pi_\omega$ if it is \emph{not} unitarily implemented there, i.e.\ if $\omega\circ\alpha_g$ is not in the folium of $\omega$.
\end{definition}

\emph{Example:} In the $m=+m_0$ ground state of an infinite ferromagnet the spin flip maps it to the $m=-m_0$ state, which lies in a different folium (just as for $\psi_\theta$ vs $\psi_{\pi/2-\theta}$): the flip has no implementer on $\HH_+$. The symmetry of the algebra is not a symmetry of the Hilbert space; the ``vacua'' are degenerate yet superselected: an observer cannot superpose them, since no local observable connects them (this is the origin of ``cat'' states being absent in infinite volume). The non-existence of the charge as an operator on $\HH_\omega$ is the algebraic form of spontaneous symmetry breaking (Goldstone's theorem, which adds the existence of gapless excitations under further assumptions, is a separate statement not derived here): the algebraic charge $Q_\Lambda=\int_\Lambda j^0$ has no weak limit on $\HH_\omega$ if the symmetry is broken, yet it acts on local operators by $[Q_\Lambda,a]\to\delta(a)$ (the infinitesimal automorphism).

\begin{keyidea}
Inequivalent representations are not pathologies; they are phases. Superselection sectors are the labels of folia, order parameters are central elements, and spontaneous symmetry breaking means the automorphism group is not unitarily implemented.
\end{keyidea}

\subsection{Thermal vs vacuum states and an outlook}

For the Weyl algebra of a free field, the thermal representation at $T>0$ is inequivalent to the vacuum representation (\cref{ch:weyl}): the thermal state has a non-zero density of particles, so $\sum_k\abs{\beta_k}^2=\infty$ in infinite volume, and by Kakutani-type reasoning the two folia are disjoint. This is the algebraic fact behind ``the Unruh/thermal state is not a state in the vacuum Fock space'', and it will return as the statement that the modular flow of the Rindler wedge is not implemented by an element of the algebra (\cref{ch:rs_bw}). In algebraic QFT, superselection sectors for charges, for example localized charged states, are the content of the DHR analysis (\cref{ch:dhr}).

\subsection*{Exercises}
\begin{exercise}\label{ex:inequiv:1}
Prove Kakutani's criterion in the easy direction: if $\sum_j(1-|\braket{\psi_j}{\psi_j'}|)<\infty$ show that $\ket{\Psi'}$ can be approximated by vectors obtained from $\ket{\Psi}$ by acting with local operators.
\end{exercise}
\begin{exercise}\label{ex:inequiv:2}
Let $\omega_{p}$ be the product state with $\rho_j=\mathrm{diag}(p,1-p)$ at every site, $p\ne1/2$. Show that $\omega_p$ and $\omega_{1/2}=\tau$ are disjoint using $M_N$. What does this say about whether a state with a finite density of excitations can be reached from the tracial state by local operations?
\end{exercise}
\begin{exercise}\label{ex:inequiv:3}
For the product state $\psi_\theta=\bigotimes_j(\cos\theta\ket\uparrow+\sin\theta\ket\downarrow)$, $\theta\in(0,\pi/4)$, show that $\psi_\theta$ is not invariant under the spin-flip automorphism $\sigma^z_j\mapsto-\sigma^z_j$ (with $\sigma^x_j\mapsto-\sigma^x_j$ so that the algebra relations hold). Show it is not unitarily implementable on $\HH_\theta$.
\end{exercise}
\begin{exercise}\label{ex:inequiv:4}
Why can a finite system never exhibit spontaneous symmetry breaking in the sense above? (Use Stone--von Neumann or the fact that $\BH$ has only one irreducible representation when $\dim\HH<\infty$.)
\end{exercise}

\section{Quasi-local algebras and the thermodynamic limit}\label{ch:quasilocal}

\begin{physmotiv}
Statistical mechanics lives in the thermodynamic limit: phase transitions, equilibrium states with infinite entropy, and the very notion of a ``phase'' exist only for infinitely many particles. In this limit $\ee^{-\beta H}$ and $\Tr$ cease to make sense, as does the Hamiltonian $H$ itself as an operator. The quasi-local algebra is the framework that survives: observables are local, states are limits of finite-volume states, and dynamics acts by automorphisms. This is also the template for AQFT, where ``finite volume'' becomes ``bounded spacetime region''.
\end{physmotiv}

\subsection{The framework}

Let $\Gamma$ be a lattice (say $\Z^d$) with a finite-dimensional algebra $\A_x=M_k(\C)$ (or CAR) at each site. For finite $\Lambda\subset\Gamma$ let $\A_\Lambda=\bigotimes_{x\in\Lambda}\A_x$. Then
\begin{itemize}
\item \textbf{isotony:} $\Lambda_1\subset\Lambda_2\Rightarrow\A_{\Lambda_1}\subset\A_{\Lambda_2}$;
\item \textbf{locality:} $[\A_{\Lambda_1},\A_{\Lambda_2}]=0$ for disjoint $\Lambda_1,\Lambda_2$ (graded for fermions);
\item the \emph{local algebra} $\A_{\rm loc}=\bigcup_\Lambda\A_\Lambda$ and the \emph{quasi-local algebra} $\A=\overline{\A_{\rm loc}}^{\norm\cdot}$.
\end{itemize}
This is a \cstar-algebra, simple, with a unique tracial state, and it is the same object regardless of the Hamiltonian. The Hamiltonian enters as an interaction $\Phi$ assigning to each finite $X\subset\Gamma$ a self-adjoint $\Phi(X)\in\A_X$, with $H_\Lambda=\sum_{X\subset\Lambda}\Phi(X)$.

\subsection{Dynamics without a Hamiltonian}

In infinite volume the formal sum $H=\sum_X\Phi(X)$ does not define an element of $\A$ (nor an operator in a general representation). Instead, define
\begin{equation}
\alpha_t(a)=\lim_{\Lambda\uparrow\Gamma}\ee^{\ii tH_\Lambda}a\,\ee^{-\ii tH_\Lambda}\qquad(a\in\A_{\rm loc}).
\end{equation}
For interactions that decay sufficiently fast (e.g.\ finite range), this limit exists in norm: this is a consequence of the \emph{Lieb--Robinson bound}, which states that the commutator $\norm{[\alpha_t(a),b]}$ for $a,b$ localized at distance $r$ is exponentially small outside the ``light cone'' $r\lesssim vt$. The limit defines a strongly continuous one-parameter group of automorphisms of $\A$ whose generator, defined on local elements, is the derivation $\delta(a)=\ii\lim[H_\Lambda,a]$.

Conclusion: \emph{dynamics is a one-parameter automorphism group $\alpha_t$, generally not inner}. There is no $H\in\A$ with $\alpha_t(a)=\ee^{\ii tH}a\ee^{-\ii tH}$. Only in a given representation $\pi_\omega$ with an $\alpha$-invariant state does $H_\omega$ exist as a self-adjoint operator (via Stone, \cref{ch:spectral}), with $H_\omega\Omega=0$ for a ground state or equilibrium reference. This is the algebraic form of the QFT lesson that the Hamiltonian depends on the vacuum representation.

\begin{whatwrong}
Different invariant states $\omega$ produce different operators $H_\omega$ on different Hilbert spaces, and the spectrum can differ (ground states have $H_\omega\ge0$; the equilibrium states at $T>0$ have $H_\omega$ unbounded in \emph{both} directions, the same pattern as the modular Hamiltonian). Never speak of ``the'' Hamiltonian of an infinite system without a representation.
\end{whatwrong}

\subsection{Thermodynamic limit of states}

Given a finite-volume Gibbs state $\omega_\Lambda=\Tr(\ee^{-\beta H_\Lambda}\,\cdot)/Z_\Lambda$, regard $\omega_\Lambda$ as a state on $\A$ (by composing with the conditional expectation onto $\A_\Lambda$, or by inserting trivial boundary conditions). The state space of $\A$ is weak-$^*$ compact (\cref{ch:states}), so a subsequence converges; its limit points are candidate infinite-volume equilibrium states. The following points matter:
\begin{enumerate}
\item The limit points are $\alpha_t$-invariant (stationary) states satisfying the \emph{KMS condition}\index{KMS condition} at inverse temperature $\beta$ (\cref{ch:kms}); this is how one defines equilibrium states directly in infinite volume, with no reference to $\ee^{-\beta H}$.
\item The set of KMS states at fixed $\beta$ is a simplex: each is a unique mixture of \emph{extremal} ones, and the extremal KMS states are the pure phases. Several extremal KMS states (e.g.\ $+$ and $-$ magnetized) at a given $\beta$ signal a first-order phase transition or symmetry breaking; a unique one means a unique phase. (The KMS states at fixed $\beta$ form a Choquet simplex, see \cite{BratteliRobinson2}, Sect.~5.3. Note that an extremal KMS state is a \emph{factor}\index{factor} state but not a \emph{pure} state of $\A$: purity would make the GNS representation irreducible, i.e.\ $\pi_\omega(\A)''=\Bnd{\HH_\omega}$ of type I, whereas for a non-tracial KMS state the factor is of type III.)
\item Extremal KMS states are \emph{factor states}: the GNS representation $\pi_\omega(\A)''$ is a factor (trivial centre, \cref{ch:factors}). Mixtures of phases have a non-trivial centre, the macroscopic order parameters (\cref{ch:inequiv}). The \emph{type} of this factor is a finer invariant: in the examples we will compute (\cref{ch:connes_class}), it is II$_1$ at $\beta=0$ and type III at generic finite temperature.
\end{enumerate}

\subsection{Why this matters for later sections}

The infinite spin chain at infinite or finite temperature is the first example of a type II$_1$ or type III algebra arising from physics: $\pi_\tau(\A)''$ is the hyperfinite II$_1$ factor, while GNS representations of product states $\omega_p$ ($p\neq1/2$) give the Araki--Woods/Powers factors of type III$_\lambda$ (\cref{ch:connes_class}). The same structure appears for a local region of a QFT: replace $\Lambda$ by a spacetime region, the lattice by a net of algebras, and the $\Z^d$ translations by Poincar\'e transformations (\cref{ch:haagkastler}).

\begin{keyidea}
The quasi-local algebra encodes kinematics (locality), dynamics is an automorphism group (not inner), equilibrium states are limits of Gibbs states characterized by the KMS condition, and phases correspond to extremal states and to the centre of the GNS von Neumann algebra.
\end{keyidea}

\subsection*{Exercises}
\begin{exercise}\label{ex:quasilocal:1}
For the non-interacting chain $H_\Lambda=\sum_{j\in\Lambda}\frac h2\sigma^z_j$, show $\alpha_t(\sigma^\pm_j)=\ee^{\pm\ii ht}\sigma^\pm_j$, so the limit exists trivially. Compute the Gibbs state limit $\omega_\beta$ and show it is a product state with $\omega_\beta(\sigma^z_j)=-\tanh(\beta h/2)$ (with $\sigma^z=\pm1$ and energy $\pm h/2$). Verify $\omega_\beta\circ\alpha_t=\omega_\beta$.
\end{exercise}
\begin{exercise}\label{ex:quasilocal:2}
Show the Hamiltonian $H=\frac h2\sum_{j\in\Z}\sigma_j^z$ is not an element of $\A$ and not the norm limit of local Hamiltonians $H_N$ (compute $\norm{H_N-H_M}$). Show that in the representation $\pi_\beta$ the generator of $\alpha_t$ exists as a self-adjoint operator but is not bounded below (consider states with local flips).
\end{exercise}
\begin{exercise}\label{ex:quasilocal:3}
Using the Ising chain with $H_\Lambda=-J\sum_j\sigma_j^z\sigma_{j+1}^z$ at $T=0$, find the two ground states (all up, all down) as product states, show they are disjoint representations, and show the infinitesimal generator $\delta$ is the same in both.
\end{exercise}
\begin{exercise}\label{ex:quasilocal:4}
State precisely why the convergence $\omega_\Lambda\to\omega$ need only be weak-$^*$ and not in norm: give an example where $\norm{\omega_\Lambda-\omega_{\Lambda'}}$ stays close to $2$ for all $\Lambda\neq\Lambda'$, e.g.\ product states with different $\theta$ restricted to the full chain.
\end{exercise}

